\documentclass[12pt,a4paper]{article}
\usepackage[utf8]{inputenc}
\usepackage[T1]{fontenc}
\usepackage[brazil]{babel}
\usepackage{amsmath,amssymb}
\usepackage{geometry}
\geometry{margin=2.5cm}
\title{Material de Estudo: M\'etodo de Runge-Kutta}
\author{Princ\'ipio de Modelagem Matem\'atica}
\date{Unidade 26}
\begin{document}
\maketitle

\section*{Objetivo}
Compreender o m\'etodo de Runge-Kutta de quarta ordem (RK4), tabelas de Butcher, e aplicar ao sistema de Lorenz.

\section*{Motiva\c{c}\~ao}

O m\'etodo de Euler tem erro $O(\Delta t)$. Queremos maior precis\~ao sem usar passos muito pequenos. A ideia do RK \'e calcular v\'arias estimativas de $f$ dentro do intervalo $[t_n, t_{n+1}]$ e combin\'a-las.

\section*{Runge-Kutta de 4\textordfeminine\ Ordem (RK4)}

Para $dy/dt = f(y, t)$:
\[
k_1 = f(y_n, t_n),
\]
\[
k_2 = f\!\left(y_n + \frac{\Delta t}{2}k_1,\; t_n + \frac{\Delta t}{2}\right),
\]
\[
k_3 = f\!\left(y_n + \frac{\Delta t}{2}k_2,\; t_n + \frac{\Delta t}{2}\right),
\]
\[
k_4 = f(y_n + \Delta t\,k_3,\; t_n + \Delta t),
\]
\[
y_{n+1} = y_n + \frac{\Delta t}{6}(k_1 + 2k_2 + 2k_3 + k_4).
\]
\begin{itemize}
  \item Erro local: $O(\Delta t^5)$; erro global: $O(\Delta t^4)$.
  \item 4 avalia\c{c}\~oes de $f$ por passo.
  \item Expl\'icito: n\~ao requer resolu\c{c}\~ao de sistema linear.
\end{itemize}

\section*{Tabela de Butcher}

A tabela de Butcher organiza os coeficientes de m\'etodos RK:
\[
\begin{array}{c|cccc}
0 & & & & \\
1/2 & 1/2 & & & \\
1/2 & 0 & 1/2 & & \\
1 & 0 & 0 & 1 & \\
\hline
& 1/6 & 2/6 & 2/6 & 1/6
\end{array}
\]

\section*{Sistema de Lorenz}

\[
\dot x = \sigma(y-x), \quad \dot y = x(\rho-z)-y, \quad \dot z = xy-\beta z.
\]
Par\^ametros cl\'assicos: $\sigma=10$, $\rho=28$, $\beta=8/3$.
\begin{itemize}
  \item Comportamento ca\'otico: trajet\'orias pr\'oximas divergem exponencialmente.
  \item Atrator estranho: dimens\~ao fractal $\approx 2{,}05$.
  \item Sensibilidade a condi\c{c}\~oes iniciais: efeito borboleta.
\end{itemize}

\section*{Exemplos Resolvidos}

\subsection*{Exemplo 1 --- RK4 para $dy/dt = y$, $y(0)=1$, $\Delta t=0{,}5$}
$k_1 = 1$; $k_2 = f(1{,}25, 0{,}25) = 1{,}25$; $k_3 = 1{,}3125$; $k_4 = f(1{,}65625, 0{,}5) = 1{,}65625$.
$y_1 = 1 + (0{,}5/6)(1 + 2{,}5 + 2{,}625 + 1{,}65625) = 1 + (0{,}5/6)\times7{,}78125 = 1{,}648$.
Exato: $e^{0{,}5} = 1{,}6487$. Erro: $0{,}04\%$.

\section*{Exerc\'icios}

\begin{enumerate}
  \item \textbf{(RK4: manual)} Aplique RK4 com $\Delta t=0{,}5$ para $dy/dt = -2y$, $y(0)=1$. Calcule $y(0{,}5)$ e compare com $e^{-1}$.
  
  \item \textbf{(Compara\c{c}\~ao)} Para $dy/dt = y^2$, $y(0)=1$ (solu\c{c}\~ao exata: $y=1/(1-t)$, explode em $t=1$): compare Euler expl\'icito e RK4 com $\Delta t=0{,}1$ at\'e $t=0{,}9$. Qual \'e mais acurado?
  
  \item \textbf{(Sistema 2D)} Resolva o sistema pend\^ulo simples:
  \[
  \dot\theta = \omega, \quad \dot\omega = -\frac{g}{\ell}\sin\theta,
  \]
  com $g=9{,}8$, $\ell=1$, $\theta(0)=0{,}5$, $\omega(0)=0$. Use RK4 com $\Delta t=0{,}1$ por 5 passos.
  
  \item \textbf{(Lorenz)} Para o sistema de Lorenz com par\^ametros cl\'assicos e $(x_0, y_0, z_0) = (1, 1, 1)$: escreva as 4 etapas do RK4 para o primeiro passo com $\Delta t=0{,}01$.
  
  \item \textbf{(Sensibilidade)} Repita o Lorenz com $(x_0, y_0, z_0) = (1{,}0001, 1, 1)$. Ap\'os quantos passos as trajet\'orias divergem visivelmente? Relate qualitat\'ivamente.
  
  \item \textbf{(Passo adapt\'ativo)} O RK-Fehlberg usa RK4 e RK5 juntos para estimar o erro e adaptar $\Delta t$. Por que isso \'e vantajoso? Em que tipo de problema isso \'e especialmente \'util?
  
  \item \textbf{(Lorenz: atrator)} O atrator de Lorenz tem dimens\~ao fractal $d_f \approx 2{,}05$. Compare com um atrator peri\'odico (ciclo limite, $d_f=1$) e um ponto fixo ($d_f=0$). O que a dimens\~ao fracion\'aria significa?
\end{enumerate}

\section*{Resumo R\'apido}
\begin{itemize}
  \item RK4: 4 estimativas de $f$; erro $O(\Delta t^4)$; 4$\times$ mais trabalho que Euler.
  \item Lorenz: $\sigma=10$, $\rho=28$, $\beta=8/3$; ca\'otico e sens\'ivel a CI.
  \item Tabela de Butcher: organiza coeficientes de qualquer m\'etodo RK.
\end{itemize}

\section*{Refer\^encias}
\begin{itemize}
  \item Lorenz, E.\ N.\ \textit{J.\ Atmos.\ Sci.}, 20, 130--141, 1963.
  \item Numerical Recipes: \textit{Chapter 17 --- Integration of ODEs}. Cambridge, 2007.
  \item Notas de aula da disciplina.
\end{itemize}

\end{document}
